\begin{proposicion}[Error en el último tablero]\label{terminal-error}
Existe una sucesión $\theta_K\to0$, dependiente sólo de las cotas exteriores, tal que en
$F\cap D_K^c$
\[
 |\log Q(B^{(K)})-L(B^{(K)})|\le\theta_K .
\]
Si $G=F\cap\{Q(B^{(K)})>0\}$, entonces $\nu(G)\ge\nu(F)-d_{n,K}$. Para $L_K^{\max}=C K\log(K+1)$ y
$\nu(G)>0$ se cumplen
\[
 \begin{aligned}
 \E_\mu\!\left[|g(B^{(K)})|\mathbf1_F\right]
 &\le\theta_K+(L_K^{\max}+\log K!)\, CK^4/n,\\
 \E_{\nu(\cdot\mid G)}[g(B^{(K)})^-]
 &\le\theta_K+L_K^{\max}\,d_{n,K}/\nu(G).
 \end{aligned}
\]
Aquí $u^-=\max\{-u,0\}$.
\end{proposicion}
\begin{proof}
En $F\cap D_K^c$ tomemos $P=(p_{ij})$ igual al punto de Bethe, extendido por cero fuera de $A$. El
Lema~\ref{terminal-pivotes} da $H_B=H_T$, $L=H_T+\log\Gamma_2$ y $Q(B^{(K)})=\per A>0$. La estacionariedad
de torres dice que $W_{ij}=a_i b_j$ en todo el soporte de $A$. Sea
$\mathcal K_T=\prod_i a_i\prod_j b_j$.
Al multiplicar los pesos de una permutación y sumar obtenemos $\per W=\mathcal K_T\,\per A$. Al
ponderar con $P_{ij}$ los logaritmos de $P_{ij}(1-P_{ij})=a_i b_j$ sobre las casillas disponibles y
usar que las masas de fila y de columna valen uno, se obtiene
\[
 H_T(P)=-\log\mathcal K_T+\sum_{ij}\log(1-P_{ij}).
\]
Por consiguiente
\[
 \log\per A-H_T(P)
 =\log\frac{\per W}{\prod_{ij}(1-P_{ij})}.
\]
El Lema~\ref{terminal-normalizador} da la primera conclusión por su uniformidad sobre los posibles
tableros en $F$. En particular $F\cap D_K^c\subseteq G$, de donde $\nu(G)\ge\nu(F)-d_{n,K}$ por el
Lema~\ref{terminal-geometria}.

En el resto de $F$ usamos directamente la entropía del predictor completo. En cada fila, $-\sum_j
p_{ij}\log p_{ij}\le\log K$, mientras que $0\le -(1-y)\log(1-y)\le y$ para $0\le y\le1$. Sumando sobre
casillas y sobre diagonales, con $\sum_e p_e=K$ y $\sum_f\ell_f=2K$, se deduce
\[
 -3K\le H_B(p)\le K\log K+2K .
\]
El Lema~\ref{coste-estatico} da entonces $|L(B^{(K)})|\le L_K^{\max}$ para $K$ suficientemente grande.
En una historia que se puede completar, $1\le Q(B^{(K)})\le K!$, de modo que
$|g(B^{(K)})|\le L_K^{\max}+\log K!$ y $g(B^{(K)})^-\le L_K^{\max}$.

Bajo $\mu$ las historias siempre se pueden completar. Aplicamos la primera conclusión en $F\cap D_K^c$
y la última cota absoluta en $F\cap D_K$, cuya probabilidad es a lo sumo $CK^4/n$ al promediar el
Lema~\ref{terminal-geometria} sobre $\pi$. Bajo $\nu(\cdot\mid G)$ la historia también es completable,
y la probabilidad de $D_K$ es a lo sumo $\nu(F\cap D_K)/\nu(G)\le d_{n,K}/\nu(G)$. La primera
conclusión en $D_K^c$ y la cota de $g^-$ en $D_K$ dan la segunda desigualdad.
\end{proof}
